\documentclass{article}
\usepackage{graphicx}
\makeatletter
\@ifundefined{rule}{\newcommand{\rule}[2]{[rule #1 by #2]}}{}
\makeatother
\title{Tables And Figures}
\author{Ferritex Bench}
\begin{document}
\maketitle
\tableofcontents

\section{Summary}
The first section introduces a schedule table and two rule-based figures that share the same narrative thread.

\begin{table}[h]
\caption{Delivery schedule}
\label{tab:schedule}
\begin{tabular}{|l|c|c|}
\hline
Track & Items & Ready \\
\hline
Layout & 4 & yes \\
Navigation & 3 & yes \\
Math & 5 & no \\
\hline
\end{tabular}
\end{table}

\begin{figure}[h]
\begin{center}
\rule{4.0in}{1.6in}
\end{center}
\caption{Pipeline sketch}
\label{fig:pipeline}
\end{figure}

\section{Comparison}
Table \ref{tab:schedule} and Figure \ref{fig:pipeline} are referenced together so the fixture mixes float metadata with normal prose.

\begin{figure}[h]
\begin{center}
\rule{3.0in}{1.4in}
\end{center}
\caption{Secondary placeholder}
\label{fig:secondary}
\end{figure}

\begin{table}[h]
\caption{Observed ratios}
\label{tab:ratios}
\begin{tabular}{|l|c|}
\hline
Metric & Value \\
\hline
Pages per section & 2.0 \\
Tables per run & 2 \\
Figures per run & 2 \\
\hline
\end{tabular}
\end{table}

Figure \ref{fig:secondary} complements Table \ref{tab:ratios}. Together they provide a compact float-heavy integration test without any external image files.
\end{document}
